\documentclass{article}

\usepackage{hyperref}
\usepackage{biblatex}

\addbibresource{main.bib}

\title{Data Aware Load Balancing}
\author{Zachary Snyder, Hamed Gorjiara, and Seungmok Lee}
\date{\today}

\begin{document}

\maketitle

\section*{Introduction}

Here, we explore load balancing in the context of information about the data
dependencies between tasks in a distributed application.  Although it is
possible to perform load balancing oblivious to the structure of the
applications, some techniques try to use more information about the application
structures to achieve better results.  We focus on dynamic load balancing
algorithms, as these seem to best fit the kinds of tasks that one encounters in
practice at present.

\section*{Background}

All of the techniques presented here are based upon work stealing.  Cilk
\autocite{RDBCEL1999} showed that work stealing could achieve existentially
optimal guarantees about the performance of a single program on a distributed
system.

The expected execution time is bounded above by \(\frac{T_1}{P} + O
(T_\infty)\), where \(T_1\) is the time it takes to serially execute the
parallelizable portions of the program, \(P\) is the number of processors, and
\(T_\infty\) is the amount of time it would take to run the program on an
infinite number of processors (critical path time).

The space required is bounded above by \(S_1 P\), where \(S_1\) is the minimum
space required to serially execute the program, and \(P\) is the number of
processors.

The expected communication overhead (measured in bytes) is bounded above by
\(O(P T_\infty (1 + n_d) S_{max} )\), where \(P\) is the number of processors,
\(T_\infty\) is the critical path time as above, \(n_d\) is the maximum number
of times that a thread synchronizes with its parent, and \(S_{max}\) is the size
of the largest activation record of any thread.

As it turns out, the parameter \(S_{max}\) ends up being the parameter of
interest, as it can be quite large.  Although these bounds are existentially
optimal, some parameters can be tweaked, and choosing the tasks to steal more
carefully than randomly may reduce some of the constants involved in the
communication overhead significantly.


\section*{State of the Art}

There are a few systems currently in use that take advantage of information
about the application structure.  We present them here to show what kind of
information is being used, and what kinds of optimizations can be made with this
information.

One of the major criticisms of Cilk is that it does not respect data locality
\autocite{KWKQISXZTLMLIR2016}.  Since transferring data is expensive when the
data is large, migrating some tasks may be more expensive than others, and may
even be more expensive than simply waiting.  This is especially true when tasks
are short-lived.  The techniques below primarily are trying to address this
concern, while building upon the work-stealing technique.


\subsection*{Hadoop Fair Scheduling}

A buzzword of a few years ago, Hadoop \autocite {JDSG2008} provides a map-reduce
model of computation to application programmers.  What makes Hadoop interesting
in the context of structure aware scheduling is that Hadoop targets a very
specific model of computation and a relatively narrow class of problems.

One of the constraints for a distributed system is fairness.  Different
processes from different users should all make progress are comparable rates.
This can be tricky to achieve when data locality is another significant
constraint, as it may be undesirable to migrate tasks, potentially compromising
fairness.

In Hadoop Fair Scheduling (HFS) \autocite{MZKEDBSSJSSIS2010} new jobs are sorted
in a list according to a hierarchical scheduling policy.  Then, each job that
will be assigned to the node that contains the input data.  If that node is busy
with other tasks, the new job can wait for a little amount of time and the
scheduler can schedule other new jobs.  With this algorithm, it is possible to
use local data nearly 100\% of the time, while maintaining relaxed a fairness
guarantee.

Spark \autocite{MZMCMJFSSIS2010} is another utility that uses HFS
\autocite{MZMCTDADJMMMMJFSSIS}, although it provides a much more general model
of computation.  %...

The reason that Hadoop and Spark get away with using HFS is that the tasks that
are run are relatively small and short-lived, and the data associated with them
is most always large.  It is less expensive to wait than to move data.  As it
happens, some believe that this is likely to become more true of all distributed
applications as computational needs grow.


\subsection*{Data Aware Work Stealing}

Waiting on large data is a good way of avoiding communication overhead, but in a
generalized and large scale environment, not all data may be large.  For such a
scenario, one might want to use something called Data Aware Work Stealing (DAWS)
\autocite{KWKQISXZTLMLIR2016}.

The idea behind DAWS is that some tasks need a lot of data and shouldn't be
migrated, but there are others that don't.  The scheduling algorithm puts two
ready queues on each node: one that is local to the machine, and one that can be
stolen from.  When a task needs to be enqueued, it is enqueued on the
appropriate queue on a node that has the task's required data.

This strategy was implemented on MATRIX \autocite{ARKWIR2013}, which is another
generalized task execution framework, this time specifically targeting machines
of epic proportions.  They also make the assumption that tasks are relatively
short.


\section*{Potential Future Work}

It seems to us that there is an issue with the model used by the authors of Cilk
\autocite{RDBCEL1999}.  The communication overhead is tallied separately from
the running time, but the reality is that the communication overhead fairly
directly contributes to the running time of any given application.  It would
seem that a better model of cost could be developed where the communications are
regarded as additional tasks in the application (data transfer has a time cost).
These tasks would be contributed, in part, by the load balancer itself.  Such a
model should more directly model the concerns that the recent work has with data
locality and the expense of migration.  It may be that a truly optimal load
balancer can rearrange the cost between the communication overhead and the
actual task scheduling to achieve a better analytical bound (if only by a large
constant) than is achieved with a data oblivious work stealing algorithm, as
expressed in the new model.


\printbibliography

\end{document}
